\documentclass[11pt,a4paper]{article}

% Csomagok
\usepackage[utf8]{inputenc}
\usepackage[magyar]{babel}
\usepackage[margin=2cm]{geometry}
\usepackage{graphicx}
\usepackage{booktabs}
\usepackage{longtable}
\usepackage{xcolor}
\usepackage{colortbl}
\usepackage{tikz}
\usepackage{pgfplots}
\usepackage{fancyhdr}
\usepackage{multicol}
\usepackage{tcolorbox}
\usepackage{fontawesome5}

\pgfplotsset{compat=1.17}

% Színek definiálása
\definecolor{clgreen}{HTML}{66CC66}
\definecolor{clblue}{HTML}{6699FF}
\definecolor{clred}{HTML}{FF6666}
\definecolor{clyellow}{HTML}{FFCC66}
\definecolor{clgray}{HTML}{CCCCCC}

% Fejléc/lábléc
\pagestyle{fancy}
\fancyhf{}
\lhead{\VAR{league_name} - \VAR{round_number}. forduló}
\rhead{\VAR{season}}
\cfoot{\thepage}

% Címoldal adatok
\title{
    \Huge \textbf{\VAR{league_name}} \\
    \Large \VAR{season} Szezon \\
    \large \VAR{round_number}. Forduló Összefoglaló
}
\author{}
\date{\VAR{generated_date}}

\begin{document}

% ==================== CÍMOLDAL ====================
\maketitle
\thispagestyle{empty}

\begin{center}
\Large Időszak: \VAR{date_range.start} -- \VAR{date_range.end}
\end{center}

\vfill

\begin{center}
\textit{Automatikusan generált riport}
\end{center}

\newpage

% ==================== EXECUTIVE SUMMARY ====================
\section*{Executive Summary}

\begin{tcolorbox}[colback=blue!5,colframe=blue!40,title=Főbb Insights]
\begin{itemize}
%% for insight in insights
    \item \VAR{insight}
%% endfor
\end{itemize}
\end{tcolorbox}

\vspace{1em}

% ==================== TABELLA & MOMENTUM ====================
\section*{Aktuális Tabella}

\begin{longtable}{clccccccl}
\toprule
\textbf{Hely} & \textbf{Csapat} & \textbf{M} & \textbf{Gy} & \textbf{D} & \textbf{V} & \textbf{Pont} & \textbf{GK} & \textbf{Forma} \\
\midrule
\endhead

%% for team in teams_metrics|sort(attribute='position')
    %% if team.position <= zone_limits.champions_league
        \rowcolor{\VAR{colors.champions_league}}
    %% elif team.position == zone_limits.europa_league
        \rowcolor{\VAR{colors.europa_league}}
    %% elif team.position == zone_limits.conference_league
        \rowcolor{\VAR{colors.conference_league}}
    %% elif team.position >= zone_limits.relegation_start
        \rowcolor{\VAR{colors.relegation}}
    %% endif
    
    \VAR{team.position} & 
    \VAR{team.team_name} & 
    \VAR{team.matches_played} & 
    \VAR{team.wins} & 
    \VAR{team.draws} & 
    \VAR{team.losses} & 
    \textbf{\VAR{team.points}} & 
    \VAR{team.goal_difference} & 
    \VAR{team.last_5_results} \\
%% endfor

\bottomrule
\end{longtable}

\begin{small}
\textcolor{clgreen}{\faSquare} Bajnokok Ligája \quad
\textcolor{clblue}{\faSquare} Európa Liga \quad
\textcolor{cyan}{\faSquare} Konferencia Liga \quad
\textcolor{clred}{\faSquare} Kiesőzóna
\end{small}

\vspace{1em}

\subsection*{Momentum - Pozícióváltozások}

\begin{center}
\begin{tikzpicture}
\begin{axis}[
    width=12cm,
    height=6cm,
    xlabel={Fordulók},
    ylabel={Helyezés},
    ymin=0, ymax=21,
    ytick={1,5,10,15,20},
    legend pos=outer north east,
    grid=major
]

%% for team in teams_metrics[:5]
    \addplot coordinates {
    %% for pos in team.position_history
        (\VAR{loop.index}, \VAR{pos})
    %% endfor
    };
    \addlegendentry{\VAR{team.team_name}}
%% endfor

\end{axis}
\end{tikzpicture}
\end{center}

\newpage

% ==================== FORDULÓ EREDMÉNYEK ====================
\section*{\VAR{round_number}. Forduló Eredményei}

\begin{center}
\begin{tabular}{lclccc}
\toprule
\textbf{Hazai} & & \textbf{Vendég} & \textbf{Eredmény} & \textbf{xG} & \textbf{Típus} \\
\midrule

%% for match in matches_enhanced
    \VAR{match.home_team} & vs & \VAR{match.away_team} & 
    \VAR{match.home_goals}:\VAR{match.away_goals} & 
    \VAR{match.home_xg|round(1)} - \VAR{match.away_xg|round(1)} &
    %% if match.match_type == 'surprise'
        \textcolor{clyellow}{\faStar}
    %% else
        --
    %% endif
    \\
%% endfor

\bottomrule
\end{tabular}
\end{center}

\vspace{1em}

\begin{tcolorbox}[colback=yellow!10,colframe=yellow!50,title=\faStar~ Meglepetés Eredmények]
%% for match in matches_enhanced
    %% if match.match_type == 'surprise'
        \textbf{\VAR{match.home_team} \VAR{match.home_goals}:\VAR{match.away_goals} \VAR{match.away_team}} -- 
        Az xG alapján (\VAR{match.home_xg|round(1)} vs \VAR{match.away_xg|round(1)}) 
        eltérő eredmény volt várható. \\[0.5em]
    %% endif
%% endfor
\end{tcolorbox}

\newpage

% ==================== STATISZTIKAI TOPLISTÁK ====================
\section*{Statisztikai Toplista}

\begin{multicols}{2}

\subsection*{\faChartLine~ Top 5 Támadás}
\begin{tabular}{lc}
\toprule
\textbf{Csapat} & \textbf{xG} \\
\midrule
%% for team in top_stats.top_attack
    \VAR{team.team_name} & \VAR{team.xg_total|round(1)} \\
%% endfor
\bottomrule
\end{tabular}

\vspace{1.5em}

\subsection*{\faChartLine~ Top 5 Védelem}
\begin{tabular}{lc}
\toprule
\textbf{Csapat} & \textbf{xGA} \\
\midrule
%% for team in top_stats.top_defense
    \VAR{team.team_name} & \VAR{team.xga_total|round(1)} \\
%% endfor
\bottomrule
\end{tabular}

\columnbreak

\subsection*{\faFire~ Legjobb Forma}
\begin{tabular}{lc}
\toprule
\textbf{Csapat} & \textbf{Index} \\
\midrule
%% for team in top_stats.best_form
    \VAR{team.team_name} & \VAR{team.form_index|round(2)} \\
%% endfor
\bottomrule
\end{tabular}

\vspace{1.5em}

\subsection*{\faArrowDown~ Legrosszabb Forma}
\begin{tabular}{lc}
\toprule
\textbf{Csapat} & \textbf{Index} \\
\midrule
%% for team in top_stats.worst_form
    \VAR{team.team_name} & \VAR{team.form_index|round(2)} \\
%% endfor
\bottomrule
\end{tabular}

\end{multicols}

\vspace{1em}

\subsection*{Teljesítmény-rés (Pont vs xPoints)}

\begin{center}
\begin{tikzpicture}
\begin{axis}[
    width=12cm,
    height=8cm,
    xlabel={Várható Pontok (xPoints)},
    ylabel={Tényleges Pontok},
    grid=major,
    legend pos=north west
]

% Referencia vonal (ideális teljesítmény)
\addplot[dashed, thick, gray] coordinates {(0,0) (30,30)};
\addlegendentry{Ideális}

% Csapatok
%% for team in teams_metrics
    \addplot[only marks, mark=*, mark size=2pt] 
    coordinates {(\VAR{team.xPoints}, \VAR{team.points})};
%% endfor

\end{axis}
\end{tikzpicture}
\end{center}

\newpage

% ==================== TRENDANALÍZIS ====================
\section*{Trendanalízis \& Figyelmeztetések}

\subsection*{\faArrowUp~ Felfelé Ívelő Csapatok}

\begin{tcolorbox}[colback=green!5,colframe=green!40]
%% for team in teams_metrics
    %% if team.category == 'upward'
        \textbf{\VAR{team.team_name}} (\VAR{team.position}. hely) -- 
        Forma: \VAR{team.form_index|round(2)}, 
        Momentum: \VAR{team.momentum} pozíció változás \\[0.3em]
    %% endif
%% endfor
\end{tcolorbox}

\subsection*{\faExclamationTriangle~ Veszélyzóna}

\begin{tcolorbox}[colback=red!5,colframe=red!40]
%% for team in teams_metrics
    %% if team.category == 'danger'
        \textbf{\VAR{team.team_name}} (\VAR{team.position}. hely) -- 
        Forma: \VAR{team.form_index|round(2)}, 
        Utolsó 5 meccs: \VAR{team.last_5_results} \\
        \textit{Sürgős beavatkozást igényelhet!} \\[0.5em]
    %% endif
%% endfor
\end{tcolorbox}

\subsection*{Alul/Felülteljesítők}

\begin{multicols}{2}

\textbf{Felülteljesítők} (szerencsések):
\begin{itemize}
%% for team in top_stats.overperformers[:3]
    \item \VAR{team.team_name}: +\VAR{team.performance_gap|round(1)} pont
%% endfor
\end{itemize}

\columnbreak

\textbf{Alulteljesítők} (szerencsétlenek):
\begin{itemize}
%% for team in top_stats.underperformers[:3]
    \item \VAR{team.team_name}: \VAR{team.performance_gap|round(1)} pont
%% endfor
\end{itemize}

\end{multicols}

\newpage

% ==================== KÖVETKEZŐ FORDULÓ ====================
\section*{Következő Forduló Előzetes}

\subsection*{\faCalendar~ Top Mérkőzések}

%% for fixture in next_round_fixtures[:3]
    \begin{tcolorbox}[colback=blue!5,colframe=blue!30]
    \textbf{\VAR{fixture.home_team} vs \VAR{fixture.away_team}} \\
    \textit{Dátum:} \VAR{fixture.date} \\
    \textit{Nehézség:} \VAR{fixture.difficulty_rating}/10 \\
    \textit{Kontextus:} \VAR{fixture.description}
    \end{tcolorbox}
    \vspace{0.5em}
%% endfor

\vfill

\begin{center}
\rule{0.5\textwidth}{0.4pt} \\
\textit{Generálva: \VAR{generated_date}} \\
\textit{Automatizált Liga-Összefoglaló Rendszer}
\end{center}

\end{document}